\documentclass[11pt,a4paper]{article}
\usepackage[margin=2.6cm]{geometry}
\usepackage{amsmath,amssymb,amsfonts}
\usepackage{graphicx}
\usepackage{booktabs}
\usepackage{xcolor}
\definecolor{navyish}{rgb}{0.10,0.15,0.45}
\usepackage[colorlinks=true,linkcolor=navyish,citecolor=navyish,urlcolor=navyish]{hyperref}
\usepackage{natbib}
\usepackage{microtype}
\setlength{\parskip}{2pt}

\title{\bfseries Dark Energy as Inverse Holographic Information:\\
A Parameter-Free $w(z)$ from an Informational Reading of the Cosmological Constant}

\author{Naser Ahani\\
\small Department of Physics, Amirkabir University of Technology, Tehran, Iran\\
\small \texttt{naser.ahani@aut.ac.ir}}

\date{July 2026}

\begin{document}
\maketitle

\begin{abstract}
\noindent
The cosmological constant problem is usually stated as a mismatch of some
122 orders of magnitude between the Planck-scale vacuum energy density and the
observed dark-energy density. Following a line of thought initiated by Banks
and developed by Padmanabhan, we first exhibit the exact algebraic identity
$\rho_\Lambda = \tfrac{3}{8}\,\rho_p/N_H$, where
$N_H \equiv A_H/4\ell_p^2 \simeq 3.3\times10^{122}$ is the
Bekenstein--Hawking information content of the de~Sitter horizon associated
with $\Lambda$. As it stands, this identity is a reformulation, not an
explanation: it relocates the question ``why is $\rho_\Lambda$ small?'' to
``why does the accessible universe carry $\sim 10^{122}$ holographic degrees
of freedom?''. The contribution of this paper is to \emph{promote} the
identity from a kinematic statement about the asymptotic horizon to a
dynamical law referred to the instantaneous future event horizon. The
resulting model is shown to coincide exactly with Li's holographic dark
energy, but with the usually free parameter fixed to $c=1$ by the
informational normalization (one bit per $4\ell_p^2$). The model therefore
makes a genuinely parameter-free prediction for the dark-energy equation of
state: $w_0 = -\tfrac13 - \tfrac23\sqrt{\Omega_{\Lambda 0}} \simeq -0.885$,
within $1\sigma$ of the DESI~DR2+CMB+PantheonPlus measurement
$w_0 = -0.838\pm0.055$, together with a definite thawing history
$w(z)\to -\tfrac13$ at high redshift and an asymptotic de~Sitter future.
The predicted sign of the evolution, however, is opposite to the central
CPL fit of DESI ($w_a>0$ here versus $w_a<0$ there), which makes the model
sharply falsifiable by forthcoming data. We close with a methodological
analysis of what is, and what is not, explained by informational
reformulations of $\Lambda$, and of the epistemic step that converts an
accounting identity into empirical physics.
\end{abstract}

\vspace{1ex}
\noindent\textbf{Keywords:} cosmological constant problem; holographic principle;
dark energy; equation of state; DESI; philosophy of cosmology.

\section{Introduction}
\label{sec:intro}

The cosmological constant problem is the observation that the vacuum energy
density suggested by quantum field theory with a Planck-scale cutoff exceeds
the observed dark-energy density by roughly 122 orders of magnitude
\citep{Weinberg1989,Martin2012}. Within $\Lambda$CDM this mismatch is treated
as an extreme cancellation to be explained by unknown ultraviolet physics.

Two developments motivate revisiting the problem from an informational
standpoint. First, the holographic principle
\citep{Bekenstein1973,Hawking1975,tHooft1993,Susskind1995} teaches that the
information content of a causal region is measured by its boundary area in
Planck units, $S = A/4\ell_p^2$; applied to cosmological horizons this
assigns the observable universe a finite information budget of order
$10^{122}$--$10^{123}$ \citep{Lloyd2002,EganLineweaver2010}. Banks
\citep{Banks2000} proposed long ago that the logical order be inverted:
$\Lambda$ is not a small energy to be explained but an \emph{input} that
fixes the dimension of the Hilbert space of an asymptotically de~Sitter
universe; Padmanabhan and Padmanabhan \citep{CosMIn2013,Padmanabhan2012}
similarly argued that the numerical value of $\Lambda$ is controlled by the
cosmic information content (``CosMIn''). Second, and quite recently, the DESI
collaboration has reported growing evidence---at the $2.8$--$4.2\sigma$ level
depending on the supernova compilation---that dark energy is \emph{not} a
constant: the equation of state today satisfies $w_0 > -1$ and evolves with
redshift \citep{DESIDR1,DESIDR2,DESIDR2ext}. If dark energy is dynamical,
frameworks in which $\rho_\Lambda$ is tied to an evolving horizon quantity
become natural candidates, whereas a strictly constant $\Lambda$ cannot
accommodate the data.

This paper does three things. Section~\ref{sec:identity} states the
informational identity in its exact form and is deliberately candid about its
epistemic status: it is an algebraic reformulation of known relations, not a
derivation of $\Lambda$. Section~\ref{sec:promotion} performs the step that
gives the reformulation empirical content: the identity, referred at each
cosmic time to the instantaneous future event horizon, becomes a dynamical
law. We show that this law is precisely Li's holographic dark energy
\citep{CKN1999,Li2004,WangReview2017} with the normally free parameter fixed
to $c=1$, so that the entire $w(z)$ history follows with \emph{no adjustable
parameters} beyond $\Omega_{\Lambda 0}$. Section~\ref{sec:pred} derives the
predictions and confronts them with DESI~DR2. Section~\ref{sec:challenges}
collects the known difficulties of the framework, and
Section~\ref{sec:philosophy} offers a methodological analysis of the move
from interpretation to prediction. Throughout, $\ell_p^2=\hbar G/c^3$,
$\rho_p = c^5/\hbar G^2$, and $\rho_\Lambda = \Lambda c^2/8\pi G$ (mass
densities).

\section{The informational identity and its epistemic status}
\label{sec:identity}

For a de~Sitter horizon of radius $R_H = \sqrt{3/\Lambda}$, the horizon area
is $A_H = 4\pi R_H^2 = 12\pi/\Lambda$, and the Bekenstein--Hawking relation
assigns it the dimensionless information content
\begin{equation}
N_H \;\equiv\; \frac{A_H}{4\ell_p^2}
\;=\; \frac{3\pi c^3}{\Lambda \hbar G}
\;\approx\; 3.3\times 10^{122},
\label{eq:NH}
\end{equation}
where the numerical value uses $\Lambda \approx 1.09\times10^{-52}\,
\mathrm{m^{-2}}$ from current data \citep{Planck2018,DESIDR2}. This is the
familiar de~Sitter entropy \citep{GibbonsHawking1977}; it agrees in order of
magnitude with Lloyd's count of the computational capacity of the universe
($\sim 10^{120}$ elementary operations) \citep{Lloyd2002} and with the
entropy budget of Egan and Lineweaver \citep{EganLineweaver2010}.

Combining Eq.~\eqref{eq:NH} with the definitions of $\rho_p$ and
$\rho_\Lambda$ gives the exact identity
\begin{equation}
\boxed{\;\rho_\Lambda \;=\; \frac{3}{8}\,\frac{\rho_p}{N_H}\;}
\label{eq:identity}
\end{equation}
The famous factor $10^{-122}$ is thus, numerically, the inverse holographic
information content of the $\Lambda$-horizon.

Three remarks fix the epistemic status of Eq.~\eqref{eq:identity}.

\emph{(i) It is not a derivation of $\Lambda$.} Both sides are functions of
the same measured $\Lambda$; the relation is an algebraic identity within
de~Sitter geometry and has, by itself, no predictive content. Any claim to
have ``derived'' the dark-energy density from
Eq.~\eqref{eq:identity} would be circular.

\emph{(ii) It is nevertheless a nontrivial reframing.} As Banks emphasized
\citep{Banks2000}, one may read the identity right-to-left: the small number
is not a fine-tuned cancellation but the reciprocal of a large integer---the
number of holographic degrees of freedom of the accessible universe. The
explanatory burden is relocated from ultraviolet cancellation mechanisms to
the question of why the cosmic Hilbert space has the dimension it has
\citep{CosMIn2013}. Relocation is not resolution, but it can redirect
research: questions about horizon information admit dynamical treatment in a
way that questions about exact cancellation of $10^{122}$ contributions do
not.

\emph{(iii) A caveat about which horizon.} Today's universe is not exactly
de~Sitter ($\Omega_{\Lambda 0}\simeq 0.7$); the Hubble radius $c/H_0$ and
the asymptotic radius $\sqrt{3/\Lambda}$ differ by a factor
$\sqrt{\Omega_{\Lambda 0}}$. Equation~\eqref{eq:identity} strictly refers to
the \emph{final} de~Sitter horizon determined by $\Lambda$ itself, which is
the formal root of the circularity noted in (i). The next section removes
this circularity by referring the information content to a horizon defined by
the expansion history rather than by $\Lambda$.

\section{From interpretation to dynamics: the promotion postulate}
\label{sec:promotion}

The step that converts Eq.~\eqref{eq:identity} from an accounting identity
into physics is the following postulate.

\medskip
\noindent\textbf{Postulate (informational dark energy).}
\emph{At every cosmic time $t$, the dark-energy density equals
$\tfrac{3}{8}\rho_p/N_H(t)$, where $N_H(t)=A_h(t)/4\ell_p^2$ is the
Bekenstein--Hawking information content of the instantaneous future event
horizon}
\begin{equation}
R_h(t) \;=\; a(t)\int_t^{\infty}\frac{c\,dt'}{a(t')} .
\label{eq:Rh}
\end{equation}

\medskip
\noindent
Substituting $N_H(t)=\pi R_h^2/\ell_p^2$ into the postulate gives
\begin{equation}
\rho_\Lambda(t) \;=\; \frac{3}{8}\,\frac{\rho_p\,\ell_p^2}{\pi R_h^2(t)}
\;=\; \frac{3c^2}{8\pi G}\,\frac{1}{R_h^2(t)},
\label{eq:rhodyn}
\end{equation}
using $\rho_p\ell_p^2 = c^2/G$. Equation~\eqref{eq:rhodyn} is recognizable:
it is exactly the holographic dark energy (HDE) density of Cohen, Kaplan and
Nelson \citep{CKN1999} with Li's choice of the future event horizon as the
infrared cutoff \citep{Li2004},
$\rho_{\rm HDE} = 3c_{\rm H}^2 M_p^2/R_h^2$
(with $M_p^{-2}=8\pi G$, $\hbar=c=1$), \emph{with the usually free
dimensionless parameter fixed to} $c_{\rm H}=1$.

This identification cuts both ways and both directions matter. In one
direction, it is a candor requirement: the dynamics explored below is not
new---HDE has a twenty-year literature, including extensive observational
fits \citep{Li2004,HuangLi2004,WangReview2017}. In the other direction, the
informational reading adds something HDE by itself lacks: a \emph{reason} for
the normalization. In standard HDE, $c_{\rm H}$ is a free parameter fitted
to data (fits typically prefer $c_{\rm H}\lesssim 1$
\citep{HuangLi2004,WangReview2017}). Here $c_{\rm H}=1$ is forced by the
Bekenstein--Hawking count of exactly one bit per $4\ell_p^2$ of horizon
area---the same factor of $1/4$ that black-hole thermodynamics fixes
\citep{Bekenstein1973,Hawking1975}. The model therefore has \emph{zero}
adjustable parameters in the dark sector: given $\Omega_{\Lambda 0}$, the
entire expansion history is determined. From a methodological point of view
(Section~\ref{sec:philosophy}) this maximal rigidity is the model's chief
virtue, since it makes the framework sharply falsifiable rather than
flexibly accommodating.

For a spatially flat universe containing matter and the component
\eqref{eq:rhodyn}, defining $\Omega_\Lambda = \rho_\Lambda/(3M_p^2H^2)$ and
$x=\ln a$, the standard HDE derivation \citep{Li2004} yields the autonomous
equation
\begin{equation}
\frac{d\Omega_\Lambda}{dx}
= \Omega_\Lambda\,(1-\Omega_\Lambda)\left(1 + 2\sqrt{\Omega_\Lambda}\right),
\label{eq:ode}
\end{equation}
and the equation of state
\begin{equation}
w(z) \;=\; -\frac13 \;-\; \frac{2}{3}\sqrt{\Omega_\Lambda(z)} .
\label{eq:wz}
\end{equation}
A self-contained derivation of Eqs.~\eqref{eq:ode}--\eqref{eq:wz} is given in
Appendix~\ref{app:derivation}. Because Eq.~\eqref{eq:wz} confines $w$ to the
open interval $(-1,-\tfrac13)$ at all finite times, the model structurally
forbids both a constant $w=-1$ and any phantom regime.

\section{Predictions and confrontation with DESI DR2}
\label{sec:pred}

\subsection{Parameter-free numbers}

Evaluating Eq.~\eqref{eq:wz} today with $\Omega_{\Lambda 0}=0.685$
\citep{Planck2018} gives the headline prediction
\begin{equation}
w_0 \;=\; -\frac13 - \frac23\sqrt{0.685} \;=\; -0.885 .
\label{eq:w0}
\end{equation}
No parameter has been tuned to obtain this number. For comparison, the DESI
DR2 combination BAO+CMB+PantheonPlus yields $w_0 = -0.838 \pm 0.055$
\citep{DESIDR2}; the prediction \eqref{eq:w0} sits within $0.9\sigma$ of the
measured central value. The same combination with the DES-Y5 supernova
sample gives $w_0 = -0.752 \pm 0.057$, about $2.3\sigma$ from \eqref{eq:w0}.

Integrating Eq.~\eqref{eq:ode} from $\Omega_{\Lambda0}=0.685$ gives the full
history (Table~\ref{tab:wz}, Fig.~\ref{fig:wz}). The model is a
\emph{thawing-to-de-Sitter} scenario: $w \to -\tfrac13$ at high redshift,
$w_0=-0.885$ today, and $w\to -1$, $\Omega_\Lambda \to 1$ in the far future.
There is no phantom regime and no big rip; the future is asymptotically
de~Sitter, closing the loop with the static identity of
Section~\ref{sec:identity}, which becomes exact in that limit. The early
dark-energy fraction is negligible where it must be: extrapolating the
matter-plus-dark-energy system (radiation neglected, so the number is
indicative only) gives $\Omega_\Lambda(z{=}1100) \sim 2\times10^{-4}$,
far below current CMB bounds on early dark energy, which sit at the percent
level.

\begin{table}[t]
\centering
\caption{Predicted history of the informational dark-energy model
($\Omega_{\Lambda0}=0.685$; no adjustable parameters).}
\label{tab:wz}
\begin{tabular}{ccc}
\toprule
$z$ & $\Omega_\Lambda(z)$ & $w(z)$\\
\midrule
0.00 & 0.685 & $-0.885$\\
0.25 & 0.551 & $-0.828$\\
0.50 & 0.441 & $-0.776$\\
1.00 & 0.295 & $-0.696$\\
2.00 & 0.160 & $-0.600$\\
3.00 & 0.104 & $-0.549$\\
10.0 & 0.026 & $-0.441$\\
$\infty$ & 0 & $-1/3$\\
\bottomrule
\end{tabular}
\end{table}

\begin{figure}[t]
\centering
\includegraphics[width=\textwidth]{wz_figure.pdf}
\caption{Left: equation of state $w(z)$ of the informational model
(solid), compared with the CPL best fits of DESI DR2 combined with CMB and
two supernova samples \citep{DESIDR2}, and with $\Lambda$CDM. Right: the
predicted dark-energy fraction $\Omega_\Lambda(z)$. The model has no free
parameters once $\Omega_{\Lambda0}$ is fixed.}
\label{fig:wz}
\end{figure}

\subsection{Agreement, tension, and falsifiability}
\label{sec:tension}

The confrontation with DESI must be reported symmetrically.

\emph{Points of agreement.} The model predicts, unavoidably, that (a) dark
energy is dynamical, so that $\Lambda$CDM is excluded within the framework;
(b) $w_0 > -1$ today, i.e.\ quintessence-like present behavior; and (c) a
present value $w_0=-0.885$ within $1\sigma$ of the
DESI+CMB+PantheonPlus measurement. Points (a) and (b) are precisely the two
qualitative headlines of the DESI results \citep{DESIDR1,DESIDR2}. It is
worth emphasizing how nontrivial (c) is: a formula with no tunable constants,
whose only inputs are the Bekenstein--Hawking factor $1/4$ and
$\Omega_{\Lambda 0}$, lands within one standard deviation of the measured
$w_0$.

\emph{Point of tension.} Projected onto the CPL parametrization
$w(a)=w_0+w_a(1-a)$ over $0<z<1$, the model gives an effective
$w_a \simeq +0.4$ (locally $+0.23$ at $z=0$): $w$ \emph{increases} into the
past toward $-\tfrac13$. The DESI CPL fits prefer the opposite sign,
$w_a \simeq -0.6$ to $-1$, i.e.\ a phantom regime $w<-1$ at $z\gtrsim0.5$
\citep{DESIDR2}. Taken at face value this is a $\sim4\sigma$ tension in
$w_a$ (a face-value statement only, since the DESI $w_0$--$w_a$ posteriors
are strongly correlated and a proper comparison must be made against the
joint contour). Three qualifications are in order. First, the CPL fit is a two-parameter
extrapolation, and its high-redshift behavior is dominated by the
parametrization rather than by data; model-independent reconstructions of
$w(z)$ carry large uncertainties beyond $z\sim0.7$, and the literature
cautions explicitly against over-reading CPL extrapolations
\citep{Nesseris2025,DESIDR2ext}. Second, a proper test of the present model
is a direct fit of its distance--redshift relation to the BAO and supernova
data, not a comparison of CPL coefficients; this fit (which involves no free
dark-sector parameters, hence no ability to hide) is the natural next step
and is left for dedicated work. Third, precisely because nothing can be
adjusted, the model is \emph{maximally exposed}: if forthcoming DESI and
supernova data confirm a genuine phantom crossing at $z\approx0.5$, the model
is refuted outright. We regard this exposure as a feature. A framework
born as an interpretation has become an empirical claim with a definite
kill condition.

\section{Known challenges}
\label{sec:challenges}

Honesty requires listing the difficulties this framework inherits and adds.

\emph{The future event horizon and causality.} The cutoff \eqref{eq:Rh}
refers to the entire future expansion history, so the present dark-energy
density ``knows about'' the future. This teleological flavor is a standing
objection to event-horizon HDE \citep{Cai2007,WangReview2017}. Two readings
are available: as a genuine causality violation (fatal), or as a global
self-consistency condition---the universe's history solves a boundary-value
problem rather than an initial-value problem, a reading with precedents in
action-at-a-distance electrodynamics and in global formulations of quantum
cosmology. Within the informational framework the second reading is at least
natural: horizon information is a global, not local, bookkeeping. But we
flag it clearly as the framework's deepest conceptual liability.

\emph{Why one bit per $4\ell_p^2$ of the cosmological horizon?} The
Bekenstein--Hawking normalization is established for black-hole horizons and
strongly motivated for de~Sitter space \citep{GibbonsHawking1977}; its
application to the time-dependent event horizon of an FRW universe is an
extrapolation. The Cohen--Kaplan--Nelson bound \citep{CKN1999} motivates
$\rho_\Lambda \lesssim M_p^2/L^2$ as an inequality; promoting it to an
equality with coefficient fixed by horizon thermodynamics is the physical
content of our postulate, and it could be wrong in either the coefficient or
the choice of surface. (Choosing the Hubble radius instead is known to fail
to produce acceleration \citep{Hsu2004}, which is itself evidence that the
choice of surface carries physical content.)

\emph{Relation to observationally preferred HDE.} Fits of standard HDE with
$c_{\rm H}$ free tend to prefer $c_{\rm H}<1$
\citep{HuangLi2004,WangReview2017}, which entails phantom behavior in the
future. Our normalization forbids this region. Should data ultimately
demand $c_{\rm H}<1$, the informational normalization---not merely the
model---would be falsified.

\emph{What $N_H$ counts.} $N_H$ is a capacity (a bound on accessible degrees
of freedom), not a count of particles or bits ``in use''. The framework does
not yet supply a microscopic account of what carries the horizon
information, nor of how matter degrees of freedom are subsumed in it. Here
the informational reading connects naturally to broader emergent-gravity
programs \citep{Padmanabhan2010,Verlinde2011}, but the connection remains
programmatic.

\section{Methodological analysis: what has and has not been explained}
\label{sec:philosophy}

The trajectory of this paper---identity, reinterpretation, promotion,
prediction---is itself of philosophical interest, and we make the structure
explicit.

\emph{Interpretations are cheap; rigid interpretations are not.} The bare
identity \eqref{eq:identity} has no empirical content: it cannot conflict
with any observation, because both sides co-vary with the measured
$\Lambda$. In Popperian terms it forbids nothing. What confers empirical
content is the promotion postulate of Section~\ref{sec:promotion}, which
extends the identity counterfactually---to all times, via a horizon defined
independently of $\Lambda$. The promoted theory forbids a great deal: it
forbids a constant $w=-1$ (the de~Sitter value is attained only
asymptotically), forbids phantom crossing, forbids any $w_0$ outside a
narrow band around $-0.89$, and fixes the full function $w(z)$. The epistemic
lesson generalizes: an interpretation earns scientific standing exactly when
it is rigid enough to be killed. The absence of adjustable parameters, often
presented apologetically, is on this view the model's principal
epistemological asset.

\emph{Relocation of explanatory burden.} Even if the model survives
observation, one should be precise about what has been explained. The value
of $\Lambda$ (equivalently $N_H$, equivalently $\Omega_{\Lambda0}$ at our
epoch) remains an unexplained input, exactly as in Banks's proposal
\citep{Banks2000} and in CosMIn \citep{CosMIn2013}. What the framework
explains, conditionally, is the \emph{structure} of dark energy: why its
density is tied to the horizon scale, why it evolves, and with what history.
The fine-tuning question is transformed, not dissolved: ``why is
$\rho_\Lambda$ small?'' becomes ``why is the cosmic information capacity
large?''---a question that at least admits dynamical treatment, since
horizon information grows with cosmic evolution and its magnitude is tied to
the age and content of the universe \citep{Lloyd2002,EganLineweaver2010}.

\emph{The Ptolemaic risk, faced honestly.} A skeptic may object that
identifying $10^{122}$ with a horizon-bit count is numerology---epicycles in
informational dress. The reply is the one physics has always given: the
difference between numerology and physics is a differential prediction. The
static identity of Section~\ref{sec:identity} cannot be distinguished from
coincidence; the dynamical model of Section~\ref{sec:promotion} can, and
DESI-class data are already adjudicating. Whichever way that adjudication
goes, the methodological point stands: informational reformulations of
$\Lambda$ must be pushed to the point of rigidity at which they expose
themselves to data, or they remain interpretive commentary.

\section{Conclusion}
\label{sec:conclusion}

We have separated, deliberately, three layers that are often run together in
holographic discussions of the cosmological constant: an exact but
non-predictive identity ($\rho_\Lambda = \tfrac38 \rho_p/N_H$); an
interpretive inversion due to Banks and Padmanabhan (read $\Lambda$ as fixing
the cosmic information capacity, not as a fine-tuned energy); and a
dynamical promotion (apply the identity at every epoch to the instantaneous
event horizon). The third layer is equivalent to holographic dark energy
with its free parameter frozen at $c_{\rm H}=1$ by the Bekenstein--Hawking
normalization, and therefore predicts, with no tunable constants,
$w_0=-0.885$---within $1\sigma$ of DESI DR2+CMB+PantheonPlus---together with
a definite thawing history whose past behavior ($w_a>0$) differs in sign
from the current CPL central fit and will be cleanly tested by forthcoming
data. A direct, parameter-free fit of the model's distance--redshift
relation to DESI BAO and supernova compilations is the immediate next step.

\paragraph{Acknowledgments.} The author thanks the instructors of the
philosophy-of-physics seminars at Sharif University of Technology for
discussions that shaped Section~\ref{sec:philosophy}.

\paragraph{Data and code availability.} No new observational data were
produced. The short numerical script that integrates Eq.~\eqref{eq:ode} and
generates Table~\ref{tab:wz} and Fig.~\ref{fig:wz} is available from the
author on request.

\appendix

\section{Derivation of the autonomous equation and the equation of state}
\label{app:derivation}

For completeness we derive Eqs.~\eqref{eq:ode}--\eqref{eq:wz} from the
postulate, following \citet{Li2004}. Units $\hbar=c=1$,
$M_p^{-2}=8\pi G$; flat FRW with matter ($\rho_m \propto a^{-3}$) and the
informational component $\rho_\Lambda = 3M_p^2/R_h^2$
(Eq.~\eqref{eq:rhodyn} with $c_{\rm H}=1$).

From the definition \eqref{eq:Rh} of the future event horizon,
\begin{equation}
\dot R_h \;=\; H R_h - 1 .
\label{eq:Rhdot}
\end{equation}
The Friedmann equation $3M_p^2H^2 = \rho_m + \rho_\Lambda$ and the
definition $\Omega_\Lambda \equiv \rho_\Lambda/3M_p^2H^2$ give
\begin{equation}
\Omega_\Lambda = \frac{1}{(H R_h)^2}
\qquad\Longrightarrow\qquad
H R_h = \Omega_\Lambda^{-1/2}.
\label{eq:HRh}
\end{equation}
With $x \equiv \ln a$ and $'\equiv d/dx = H^{-1}d/dt$,
Eq.~\eqref{eq:Rhdot} gives
\begin{equation}
\frac{R_h'}{R_h} \;=\; 1 - \frac{1}{HR_h} \;=\; 1 - \sqrt{\Omega_\Lambda}.
\label{eq:Rhprime}
\end{equation}
Differentiating $\rho_\Lambda = 3M_p^2/R_h^2$ and comparing with the
continuity equation $\rho_\Lambda' = -3(1+w)\rho_\Lambda$ yields the
equation of state directly:
\begin{equation}
\rho_\Lambda' = -2\rho_\Lambda\frac{R_h'}{R_h}
= -2\rho_\Lambda\left(1-\sqrt{\Omega_\Lambda}\right)
\;\;\Longrightarrow\;\;
w = -\frac13 - \frac23\sqrt{\Omega_\Lambda},
\end{equation}
which is Eq.~\eqref{eq:wz}. Since $0<\Omega_\Lambda<1$ at all finite times,
$w$ is confined to $(-1,-\tfrac13)$.

For the autonomous equation, differentiate the logarithm of
Eq.~\eqref{eq:HRh},
\[
\Omega_\Lambda'/\Omega_\Lambda = -2\left(H'/H + R_h'/R_h\right).
\]
Differentiating the Friedmann equation and using
$\rho_m' = -3\rho_m$ gives
\begin{equation}
\frac{H'}{H} = -\frac32\left(1-\Omega_\Lambda\right)
- \Omega_\Lambda\left(1-\sqrt{\Omega_\Lambda}\right).
\end{equation}
Combining the last three relations,
\begin{equation}
\frac{\Omega_\Lambda'}{\Omega_\Lambda}
= 3(1-\Omega_\Lambda) + 2\Omega_\Lambda(1-\sqrt{\Omega_\Lambda})
- 2(1-\sqrt{\Omega_\Lambda})
= (1-\Omega_\Lambda)\left(1+2\sqrt{\Omega_\Lambda}\right),
\end{equation}
which is Eq.~\eqref{eq:ode}. The fixed points are
$\Omega_\Lambda=0$ (unstable; matter domination, $w\to-\tfrac13$) and
$\Omega_\Lambda=1$ (attractor; de~Sitter, $w\to-1$), confirming the
thawing-to-de-Sitter history quoted in the text.

\begin{thebibliography}{99}
\small

\bibitem[Weinberg(1989)]{Weinberg1989}
S.~Weinberg, ``The cosmological constant problem,''
\emph{Rev. Mod. Phys.} \textbf{61}, 1 (1989).

\bibitem[Martin(2012)]{Martin2012}
J.~Martin, ``Everything you always wanted to know about the cosmological
constant problem (but were afraid to ask),''
\emph{C. R. Physique} \textbf{13}, 566 (2012), arXiv:1205.3365.

\bibitem[Bekenstein(1973)]{Bekenstein1973}
J.~D. Bekenstein, ``Black holes and entropy,''
\emph{Phys. Rev. D} \textbf{7}, 2333 (1973).

\bibitem[Hawking(1975)]{Hawking1975}
S.~W. Hawking, ``Particle creation by black holes,''
\emph{Commun. Math. Phys.} \textbf{43}, 199 (1975).

\bibitem['t Hooft(1993)]{tHooft1993}
G.~'t~Hooft, ``Dimensional reduction in quantum gravity,''
arXiv:gr-qc/9310026.

\bibitem[Susskind(1995)]{Susskind1995}
L.~Susskind, ``The world as a hologram,''
\emph{J. Math. Phys.} \textbf{36}, 6377 (1995), arXiv:hep-th/9409089.

\bibitem[Gibbons and Hawking(1977)]{GibbonsHawking1977}
G.~W. Gibbons and S.~W. Hawking, ``Cosmological event horizons,
thermodynamics, and particle creation,''
\emph{Phys. Rev. D} \textbf{15}, 2738 (1977).

\bibitem[Banks(2000)]{Banks2000}
T.~Banks, ``Cosmological breaking of supersymmetry?,''
\emph{Int. J. Mod. Phys. A} \textbf{16}, 910 (2001), arXiv:hep-th/0007146.

\bibitem[Padmanabhan and Padmanabhan(2013)]{CosMIn2013}
T.~Padmanabhan and H.~Padmanabhan, ``CosMIn: The solution to the
cosmological constant problem,''
\emph{Int. J. Mod. Phys. D} \textbf{22}, 1342001 (2013), arXiv:1302.3226.

\bibitem[Padmanabhan(2012)]{Padmanabhan2012}
T.~Padmanabhan, ``The physical principle that determines the value of the
cosmological constant,'' arXiv:1210.4174.

\bibitem[Padmanabhan(2010)]{Padmanabhan2010}
T.~Padmanabhan, ``Thermodynamical aspects of gravity: new insights,''
\emph{Rep. Prog. Phys.} \textbf{73}, 046901 (2010), arXiv:0911.5004.

\bibitem[Verlinde(2011)]{Verlinde2011}
E.~Verlinde, ``On the origin of gravity and the laws of Newton,''
\emph{JHEP} \textbf{04}, 029 (2011), arXiv:1001.0785.

\bibitem[Lloyd(2002)]{Lloyd2002}
S.~Lloyd, ``Computational capacity of the universe,''
\emph{Phys. Rev. Lett.} \textbf{88}, 237901 (2002), arXiv:quant-ph/0110141.

\bibitem[Egan and Lineweaver(2010)]{EganLineweaver2010}
C.~A. Egan and C.~H. Lineweaver, ``A larger estimate of the entropy of the
universe,'' \emph{Astrophys. J.} \textbf{710}, 1825 (2010), arXiv:0909.3983.

\bibitem[Cohen et al.(1999)]{CKN1999}
A.~G. Cohen, D.~B. Kaplan, and A.~E. Nelson, ``Effective field theory, black
holes, and the cosmological constant,''
\emph{Phys. Rev. Lett.} \textbf{82}, 4971 (1999), arXiv:hep-th/9803132.

\bibitem[Li(2004)]{Li2004}
M.~Li, ``A model of holographic dark energy,''
\emph{Phys. Lett. B} \textbf{603}, 1 (2004), arXiv:hep-th/0403127.

\bibitem[Huang and Li(2004)]{HuangLi2004}
Q.-G. Huang and M.~Li, ``The holographic dark energy in a non-flat
universe,'' \emph{JCAP} \textbf{08}, 013 (2004), arXiv:astro-ph/0404229.

\bibitem[Wang et al.(2017)]{WangReview2017}
S.~Wang, Y.~Wang, and M.~Li, ``Holographic dark energy,''
\emph{Phys. Rept.} \textbf{696}, 1 (2017), arXiv:1612.00345.

\bibitem[Cai(2007)]{Cai2007}
R.-G. Cai, ``A dark energy model characterized by the age of the universe,''
\emph{Phys. Lett. B} \textbf{657}, 228 (2007), arXiv:0707.4049.

\bibitem[Hsu(2004)]{Hsu2004}
S.~D.~H. Hsu, ``Entropy bounds and dark energy,''
\emph{Phys. Lett. B} \textbf{594}, 13 (2004), arXiv:hep-th/0403052.

\bibitem[DESI Collaboration(2024)]{DESIDR1}
DESI Collaboration, ``DESI 2024 VI: Cosmological constraints from the
measurements of baryon acoustic oscillations,''
\emph{JCAP} \textbf{02}, 021 (2025), arXiv:2404.03002.

\bibitem[DESI Collaboration(2025)]{DESIDR2}
DESI Collaboration (M.~Abdul Karim \emph{et al.}), ``DESI DR2 Results II:
Measurements of baryon acoustic oscillations and cosmological constraints,''
arXiv:2503.14738.

\bibitem[Lodha et al.(2025)]{DESIDR2ext}
K.~Lodha \emph{et al.} (DESI), ``Extended dark energy analysis using DESI DR2
BAO measurements,'' arXiv:2503.14743.

\bibitem[Nesseris et al.(2025)]{Nesseris2025}
S.~Nesseris, Y.~Akrami, and G.~D. Starkman, ``To CPL, or not to CPL? What we
have not learned about the dark energy equation of state,''
arXiv:2503.22529.

\bibitem[Planck Collaboration(2020)]{Planck2018}
Planck Collaboration, ``Planck 2018 results. VI. Cosmological parameters,''
\emph{Astron. Astrophys.} \textbf{641}, A6 (2020), arXiv:1807.06209.

\end{thebibliography}

\end{document}
